\subsection{应用}

命题 2. 设 A 是交换环，\(\mathfrak{S}\) 是 A 的理想，B 是交换 Noether A-代数，且 B 包含于 B 的 Jacobson 根中。则每个有限生成 B-模 M 都是关于 \(\mathfrak{S}\) 的理想 Hausdorff A-模。

更一般地，我们将看到，对于每个有限生成 A-模 N，\(N \otimes_{A} M\) 关于 \(\Im\)-进拓扑是 Hausdorff 的。因为 \(\mathrm{N}_{(\mathrm{B})} = \mathrm{N} \otimes_{\mathrm{A}} \mathrm{B}\) 是有限生成 B-模，而 B-模 \(N \otimes_{A} M\) 根据张量积的结合律典范地识别为 \(\mathrm{N}_{(\mathrm{B})} \otimes_{\mathrm{B}} M\)。令 L 为 B 的 Jacobson 根；由于 \(\Im B\) 包含于 L 中，\(\Im\)-进拓扑在 \(N \otimes_{A} M\) 上可识别为比 \(\Im\)-进拓扑在 \(\mathrm{N}_{(\mathrm{B})} \otimes_{\mathrm{B}} M\) 上诱导的拓扑更细；但后一个拓扑是 Hausdorff 的，因为 \(\mathrm{N}_{(\mathrm{B})} \otimes_{\mathrm{B}} M\) 是有限生成 B-模（第 1 节，例 1），结论得证。

命题 3. 设 A 是交换环，B 是交换 A-代数，\(\mathfrak{I}\) 是 A 的理想，M 是 B-模。假设 B 是 Noether 环及平坦 A-模，且 M 关于 \(\mathfrak{S}\) B 是理想 Hausdorff 的。以下条件等价：

\begin{enumerate}
\def\labelenumi{(\alph{enumi})}
\item
  M 是平坦 B-模。
\item
  \(\mathbf{M}\) 是平坦 A-模，且 \(\mathbf{M} / \Im \mathbf{M} = \mathbf{M} / (\Im \mathbf{B})\mathbf{M}\) 是平坦的 (B/3B)-模。
\end{enumerate}

若进一步地，典范同态 \(\mathbf{A} / \Im \to \mathbf{B} / \Im \mathbf{B}\) 是双射，则条件 (a) 和 (b) 还等价于：

\begin{enumerate}
\def\labelenumi{(\alph{enumi})}
\setcounter{enumi}{2}
\tightlist
\item
  M 是平坦 A-模。
\end{enumerate}

条件 (a) 蕴含 (b)，这是由第 I 章 §2 第 7 节命题 8 的推论 2 和 3，以及 M/Im 同构于 M ⊗B (B/ImB) 这一事实得到的。设条件 (b) 成立；为了证明 M 是平坦 B-模，我们将把第 2 节的定理 1 中的 A 换为 B、S 换为 SIB。因此，只需证明典范映射 \(f\colon \mathfrak{B} \otimes_{\mathbf{B}} \mathbf{M} \to \mathfrak{M}\) 是单射。令 \(f_1\) 为典范映射 \(\mathfrak{I} \otimes_{\mathbf{A}} \mathbf{B} \to \mathfrak{B}\)，令 \(f_2\) 为典范同构 \(\mathfrak{I} \otimes_{\mathbf{A}} \mathbf{M} \to (\mathfrak{I} \otimes_{\mathbf{A}} \mathbf{B}) \otimes_{\mathbf{B}} \mathbf{M}\)；易于验证，\(f \circ (f_1 \circ 1_{\mathbf{M}}) \circ f_2\) 是典范映射 \(f'\colon \mathfrak{I} \otimes_{\mathbf{A}} \mathbf{M} \to \mathfrak{IM}\)。现在 \(f'\) 是同构，因为 M 是平坦 A-模；而 \(f_1\) 是同构，因为 B 关于 A 平坦；于是 \(f\) 是同构。

令 \(\rho: \mathring{\mathrm{A}}/\mathfrak{I} \to \mathrm{B}/\mathfrak{I}\mathrm{B}\) 为典范同态；(A/ \(\mathfrak{I}\) )-模结构在 M/ \(\mathfrak{I}\) M 上经由 \(\rho\) 从其 (B/ \(\mathfrak{I}\) B)-模结构导出，所得结构同构于 M \(\otimes_{\mathrm{A}} (\mathrm{A}/\mathfrak{I})\) 上的模结构。因此，若 M 是平坦 A-模，则 M/ \(\mathfrak{I}\) M 是平坦的 (A/ \(\mathfrak{I}\) )-模，也是 (B/ \(\mathfrak{I}\) B)-模（当 \(\rho\) 为同构时）；由此我们在该情形证明了 (c) \(\Rightarrow\) (b)。

推论。设 A 是交换 Noether 环，\(\Im\) 是 A 的理想，\(\hat{\mathbf{A}}\) 是 A 关于 \(\Im\)-进拓扑的 Hausdorff 完备化，M 是一个 \(\hat{\mathbf{A}}\)-模，且关于 \(\Im\hat{\mathbf{A}}\) 是理想 Hausdorff 模。M 为平坦 A-模的充要条件是 M 为平坦 \(\hat{\mathbf{A}}\)-模。

事实上，我们知道 \(\hat{\mathbf{A}}\) 是 Noether 环（\(\S 3\)，第 4 节，命题 8）及平坦 A-模（\(\S 3\)，第 4 节，定理 3），\(\Im \hat{\mathbf{A}} = \hat{\mathfrak{I}}\)（\(\S 2\)，第 12 节，命题 16），且典范同态 \(\mathrm{A} / \Im \to \hat{\mathrm{A}} / \hat{\mathfrak{I}}\) 是双射（\(\S 2\)，第 12 节，命题 15）；因此可应用命题 3。

命题 4. 设 A、B 是两个交换 Noether 环，h: A → B 是环同态，∑ 是 A 的理想，∑ 是 B 的一个包含 ∑B 且包含于 B 的 Jacobson 根中的理想。令 A 为 A 关于 ∑-进拓扑的 Hausdorff 完备化，B 为 B 关于 ∑-进拓扑的 Hausdorff 完备化；h 关于这些拓扑连续，因而 h: A → B 使 B 成为 A-代数。设 M 是有限生成 B-模，M 是 M 关于 ∑-进拓扑的 Hausdorff 完备化；以下性质等价：

\begin{enumerate}
\def\labelenumi{(\alph{enumi})}
\item
  M 是平坦 A-模。
\item
  \(\hat{\mathbf{M}}\) 是平坦 A-模。
\item
  \(\hat{\mathbf{M}}\) 是平坦 \(\hat{\mathbf{A}}\)-模。
\end{enumerate}

由于 B 关于 \(\mathfrak{L}\)-进拓扑是 Zariski 环，\(\hat{\mathbf{B}}\) 是忠实平坦 B-模（§ 3，第 5 节，命题 9），且 \(\hat{\mathbf{M}}\) 典范同构于 \(\mathbf{M} \otimes_{\mathbf{B}} \hat{\mathbf{B}}\)（§ 3，第 4 节，定理 3）；立即可验证，这个典范同构把 \(\hat{\mathbf{M}}\) 上的 A-模结构同构到由 M 上的结构导出的 \(\mathbf{M} \otimes_{\mathbf{B}} \hat{\mathbf{B}}\) 上的 A-模结构。将第 I 章 § 3 第 2 节的命题 4 中的 R 换为 B、S 换为 A、E 换为 \(\hat{\mathbf{B}}\)、F 换为 M，可知 M 为平坦 A-模的充要条件是 \(\hat{\mathbf{M}}\) 为平坦 A-模。此外，\(\hat{\mathbf{M}}\) 是有限生成 \(\hat{\mathbf{B}}\)-模，且 \(\mathfrak{J}\hat{\mathbf{B}}\) 包含于 \(\hat{\mathfrak{L}} = \mathfrak{L}\hat{\mathbf{B}}\)，因而包含于 \(\hat{\mathbf{B}}\) 的 Jacobson 根中（§ 3，第 4 节，命题 8）；所以 \(\hat{\mathbf{M}}\) 是理想 Hausdorff \(\hat{\mathbf{A}}\)-模，且关于 \(\mathfrak{J}\hat{\mathbf{A}}\)（命题 2）。因此，根据命题 3 的推论，条件 (b) 和 (c) 等价。


